\section{The formal degree with the specified Haar measure}
\label{fd-section}

We determine the formal degree in the normalization used by the mass
formula. This also checks the meaning of the representation parameter:
the source specifies a smallest compact type, rather than writing an
infinitesimal character into its definition of a discrete series.
The calculation below uses classical discrete-series theorems, stated
explicitly. The measure comparison, scalar integral, and the application
to the source's precise definition are proved here.

For an integer $a\ge0$, write
\[
                     s_a=\prod_{j=0}^{a-1}j!,\qquad s_0=1.
\]
Let $p,q\ge1$, put $n=p+q$, and set
\[
 G=U(p,q)=\{g\in\operatorname{GL}_n(\mathbb C):g^*Jg=J\},
 \quad J=\operatorname{diag}(I_p,-I_q),\quad
 K=U(p)\times U(q).
\]
All Lebesgue measures on complex matrices below are the products
of $d(\Re z)\,d(\Im z)$ over their entries. Haar measure on a compact
group is a probability measure unless stated otherwise.

\subsection{Two elementary matrix integrals}

Put
\[
 \mathcal B_{p,q}=\{Z\in M_{p,q}(\mathbb C):I_q-Z^*Z>0\},
\]
where $>0$ means positive definite.

\begin{lemma}\label{fd-matrix-integrals}
For $\alpha>-1$ and $u>p+q-1$, respectively,
\begin{align}
 \int_{\mathcal B_{p,q}}\det(I_q-Z^*Z)^\alpha\,dZ
   &=\pi^{pq}\prod_{j=0}^{q-1}
                  \frac{\Gamma(\alpha+j+1)}
                       {\Gamma(\alpha+p+j+1)},\label{fd-ball-integral}\\
 \int_{M_{p,q}(\mathbb C)}\det(I_q+Z^*Z)^{-u}\,dZ
   &=\pi^{pq}\prod_{j=0}^{q-1}
                  \frac{\Gamma(u-p-j)}{\Gamma(u-j)}.
                                                   \label{fd-affine-integral}
\end{align}
Here $\Gamma(v)=\int_0^\infty t^{v-1}e^{-t}\,dt$ for $v>0$.
\end{lemma}
\begin{proof}
First, for a vector $z\in\mathbb C^p$, radial integration gives
\begin{align*}
 \int_{|z|<1}(1-|z|^2)^\alpha\,dz
       &=\pi^p\frac{\Gamma(\alpha+1)}
                         {\Gamma(\alpha+p+1)},\\
 \int_{\mathbb C^p}(1+|z|^2)^{-u}\,dz
       &=\pi^p\frac{\Gamma(u-p)}{\Gamma(u)}\quad(u>p).
\end{align*}
Indeed the radial volume element is
$2\pi^p t^{2p-1}dt/\Gamma(p)$. Its constant follows by
integrating $e^{-|z|^2}$ in Cartesian and polar coordinates.
The substitution $x=t^2$ reduces the two integrals to the
beta integral. The identity
$B(a,b)=\Gamma(a)\Gamma(b)/\Gamma(a+b)$ follows by setting
$r=x+y$, $v=x/(x+y)$ in the product of the two defining
gamma integrals. All integrands are nonnegative, so the changes
of variables and repeated integration follow from Tonelli's theorem.

Write a matrix as $Z=[W,z]$, with $q-1$ columns in $W$.
The Schur complement gives
\[
 \det(I_q\pm Z^*Z)
  =\det(I_{q-1}\pm W^*W)
       \bigl(1\pm z^*(I_p\pm WW^*)^{-1}z\bigr).
 \tag{FD1}
\]
For the minus sign, membership in $\mathcal B_{p,q}$ is
equivalent to $W\in\mathcal B_{p,q-1}$ and
$z^*(I_p-WW^*)^{-1}z<1$. The equivalence of
$I_{q-1}-W^*W>0$ and $I_p-WW^*>0$ follows by comparing
their nonzero singular values.
In either case substitute $z=(I_p\pm WW^*)^{1/2}v$.
The real Jacobian is $\det(I_p\pm WW^*)$, which equals
$\det(I_{q-1}\pm W^*W)$. The equality follows, for example,
by expanding $\det\bigl(\begin{smallmatrix}I&W\\
\mp W^*&I\end{smallmatrix}\bigr)$ in its two Schur complements.

If the two integrals are denoted by $I_{p,q}(\alpha)$ and
$J_{p,q}(u)$, this gives
\begin{align*}
 I_{p,q}(\alpha)
   &=\pi^p\frac{\Gamma(\alpha+1)}{\Gamma(\alpha+p+1)}
                                      I_{p,q-1}(\alpha+1),\\
 J_{p,q}(u)
   &=\pi^p\frac{\Gamma(u-p)}{\Gamma(u)}
                                      J_{p,q-1}(u-1).
\end{align*}
At $q=0$ both integrals are one. Iteration proves the formulas,
including convergence under exactly the indicated sufficient
conditions.
\end{proof}

\subsection{The compact-dual normalization}

For a block matrix
$g=\bigl(\begin{smallmatrix}A&B\\C&D\end{smallmatrix}\bigr)\in G$,
the action on $\mathcal B_{p,q}$ is
\[
                 g\cdot Z=(AZ+B)(CZ+D)^{-1}.
\]
The denominator is invertible: otherwise a nonzero vector
$(Zv,v)$ of strictly negative norm would be carried to a
vector with second component zero, whose norm is nonnegative.
The stabilizer of $0$ is $K$, and the action is transitive.
For instance, the positive Hermitian block matrix
\[
 g_Z=
 \begin{pmatrix}
 (I_p-ZZ^*)^{-1/2}&Z(I_q-Z^*Z)^{-1/2}\\
 Z^*(I_p-ZZ^*)^{-1/2}&(I_q-Z^*Z)^{-1/2}
 \end{pmatrix}
 \tag{FD2}
\]
belongs to $G$ and carries $0$ to $Z$. The needed identities
between its off-diagonal blocks follow by a singular-value
decomposition of $Z$.

\begin{lemma}\label{fd-quotient-measure}
The measure
\[
                   d\nu_0(Z)=\det(I_q-Z^*Z)^{-n}\,dZ
 \tag{FD3}
\]
is $G$-invariant. It defines a Haar measure $\mu_0$ on $G$ by
\[
       \int_G F(g)\,d\mu_0(g)
          =\int_{\mathcal B_{p,q}}\int_K F(g_Zk)\,dk\,d\nu_0(Z).
 \tag{FD4}
\]
The pushforward of Haar probability measure on $U(n)$ to
$U(n)/(U(p)\times U(q))$ has, in its graph chart
$M_{p,q}(\mathbb C)$, density
\[
             c_{p,q}\det(I_q+Z^*Z)^{-n}\,dZ,\qquad
             c_{p,q}=\frac{s_n}{\pi^{pq}s_ps_q}.
 \tag{FD5}\label{fd-cpq}
\]
\end{lemma}
\begin{proof}
Put $Z'=g\cdot Z$ and $R=CZ+D$. Directly from $g^*Jg=J$,
\[
                  I_q-Z'^*Z'=R^{-*}(I_q-Z^*Z)R^{-1}.
 \tag{FD6}
\]
The complex derivative sends a variation $H$ to
$(A-Z'C)HR^{-1}$. The block determinant identity
\[
                         \det(A-Z'C)\det R=\det g
\]
follows by eliminating the upper-right block after replacing
the second block column of $g$ by its product with
$\bigl(\begin{smallmatrix}Z\\I_q\end{smallmatrix}\bigr)$.
The complex determinant of the derivative is
$(\det g)^q\det R^{-n}$; its real Jacobian is
$|\det R|^{-2n}$, since $|\det g|=1$.
Together with (FD6) this proves invariance of (FD3).
Changing the section in (FD4) multiplies it on the right by an
element of $K$, which does not change the inner integral.
Left translation by $G$ changes the section in this way and
preserves $\nu_0$, proving the Haar assertion.

The compact quotient is the Grassmannian of $q$-planes
in $\mathbb C^{p+q}$. Its graph chart sends $Z$ to the
column space of $\bigl(\begin{smallmatrix}Z\\I_q\end{smallmatrix}\bigr)$.
The same calculation for $g\in U(n)$ has plus signs:
\[
                   I_q+Z'^*Z'=R^{-*}(I_q+Z^*Z)R^{-1}.
 \tag{FD7}
\]
It therefore determines the invariant density in the chart
up to a positive constant. Haar measure gives a smooth
positive density in local coordinates. The complement of
the graph chart is the zero set of a nonzero Pl\"ucker
coordinate, and has measure zero. To justify this last
assertion without a measure convention, cover the
Grassmannian by its finitely many graph charts; a nonzero
complex polynomial in each chart vanishes on a set of
Lebesgue measure zero, by induction on the number of
variables and Fubini's theorem.
Formula \eqref{fd-affine-integral} at $u=n$ gives total mass
\[
 \pi^{pq}\prod_{j=0}^{q-1}\frac{(q-j-1)!}{(n-j-1)!}
                        =\pi^{pq}\frac{s_ps_q}{s_n}.
\]
Normalizing the total mass to one proves (FD5).
\end{proof}

We now specify the Lie-algebra density used in the source.
For either real form $\mathfrak u(p,q)$ or $\mathfrak u(n)$
inside $M_n(\mathbb C)$, let $\omega_{\mathfrak g}$ be the
positive Lebesgue measure characterized by
\[
 (X,Y)\longmapsto X+iY
       \quad\hbox{carrying}\quad
 \omega_{\mathfrak g}\otimes\omega_{\mathfrak g}
       \quad\hbox{to ordinary Lebesgue measure on }M_n(\mathbb C).
 \tag{FD8}
\]
This map is a real-linear isomorphism. On the diagonal write
$X_{ii}=ix_i$ with $x_i\in\mathbb R$; off the diagonal
use one complex coordinate $z_{ij}$ for each $i<j$.
The other entry is $-\varepsilon_i\varepsilon_j\bar z_{ij}$,
where $\varepsilon_i$ is the sign of the Hermitian form.
For each pair $i<j$, the real determinant in (FD8) has
absolute value $4$; on each diagonal coordinate it has
absolute value $1$. Consequently
\[
 \omega_{\mathfrak g}
    =2^{n(n-1)/2}\prod_i dx_i
                  \prod_{i<j}d(\Re z_{ij})\,d(\Im z_{ij}).
 \tag{FD9}
\]
This also proves uniqueness. In particular the density is
the same in these coordinates for the two real forms.

Call a Haar measure on $G$ compact-dual normalized if its
density at the identity relative to $\omega_{\mathfrak u(p,q)}$
equals that of Haar probability measure on $U(n)$ relative
to $\omega_{\mathfrak u(n)}$. Density may equivalently be
defined by the limit of the ratio of the measures of
exponential images of shrinking coordinate balls: the
Jacobian is continuous and the derivative of the exponential
map at zero is the identity.

\begin{proposition}\label{fd-measure-match}
There is exactly one compact-dual normalized Haar measure, and it is
\[
                             \mu_{\rm cd}=c_{p,q}\mu_0.
 \tag{FD10}\label{fd-cd-measure}
\]
This is the source's measure specified by its compact
probability and equal exponential-density conditions.
\end{proposition}
\begin{proof}
For both groups the decomposition at the identity into
the block-diagonal compact algebra and its complementary
off-diagonal coordinates has
\[
          \omega_{\mathfrak g}
                =2^{pq}\omega_{\mathfrak k}\,dZ.
 \tag{FD11}
\]
Indeed (FD9) contributes one additional factor $2$ for
each of the $pq$ cross-block complex entries.
Let $b_K$ be the density of Haar probability measure on
$K$ relative to $\omega_{\mathfrak k}$.
For $G$, the local product map $(Z,k)\mapsto g_Zk$ has
derivative the direct-sum identification at $(0,1)$.
By (FD3)--(FD4), the identity density of $\mu_0$ is
$b_K/2^{pq}$.
For the compact group use the analogous section obtained
by orthonormalizing the graph of $Z$; its derivative has
off-diagonal blocks $Z,-Z^*$ and the same coordinate
Jacobian. Formula (FD5) gives identity density
$c_{p,q}b_K/2^{pq}$.
Their ratio proves (FD10). All Haar measures on a fixed
Lie group are positive scalar multiples of one another,
so the density also proves uniqueness.
\end{proof}

\subsection{Classical representation-theoretic inputs}

For a strongly continuous irreducible unitary discrete-series
representation $\pi$ of $G$, its formal degree with respect to
a Haar measure $\mu$ is the positive number $d_\mu(\pi)$ in
Schur orthogonality; in particular
\[
  \int_G|\langle\pi(g)v,w\rangle|^2\,d\mu(g)
      =d_\mu(\pi)^{-1}\|v\|^2\|w\|^2.
 \tag{FD12}
\]
It follows directly that
$d_{a\mu}(\pi)=a^{-1}d_\mu(\pi)$ for $a>0$.
A discrete series is an irreducible unitary representation
occurring as a closed subrepresentation of $L^2(G,\mu)$.

We use the following classical facts. They are background
theorems, not new deductions or unproved assertions from
the submitted archive.

\begin{enumerate}
\item Harish--Chandra's classification and formal-degree theorem.
Discrete series have regular integral Harish--Chandra parameters,
considered modulo the compact Weyl group.
For fixed $G$ and $\mu$ there is a positive constant $b_\mu$,
independent of the parameter and its chamber, such that
\[
                    d_\mu(\pi_\lambda)
                       =b_\mu\prod_{i<j}|\lambda_i-\lambda_j|.
 \tag{FD13}\label{fd-classical-hc}
\]
Here the integrality convention for $U(p,q)$ is
$\lambda_i-(n-1)/2\in\mathbb Z$.
The constant is left undetermined at this stage.
\cite[Theorem 16, \S41]{harish-chandra-ds-ii}
\item Blattner's formula: choose the positive roots making
$\lambda$ dominant, let $\rho_{\rm c}$ and $\rho_{\rm n}$
be the half sums of its compact and noncompact positive
roots, and let $P_{\rm n}(\beta)$ count the expressions of
$\beta$ as a sum of noncompact positive roots. For a
compact-dominant integral highest weight $\mu$, its
multiplicity is
\[
 m_\mu(\pi_\lambda)=
  \sum_{w\in W_K}\operatorname{sgn}(w)
       P_{\rm n}\bigl(w(\mu+\rho_{\rm c})
                                -\lambda-\rho_{\rm n}\bigr).
 \tag{FD14}\label{fd-blattner}
\]
The partition function is finite because these roots lie
in an open half-space.
\cite[III.1, theorem (2)]{duflo-heckman-vergne}
The stated highest-weight convention is also checked
against the explicit $U(p,q)$ convention in
\cite[\S2.2]{baldoni-vergne}.
\item The characterization by a compact type
\cite[Theorem 1.1]{parthasarathy-minimal}:
put $\nu=\lambda-\rho$, with $\rho=\rho_{\rm c}+\rho_{\rm n}$,
and suppose $\nu$ is dominant integral. An irreducible
unitary representation containing the compact type of
highest weight $\nu+2\rho_{\rm n}$, and containing no
compact-dominant type $\nu+2\rho_{\rm n}-\alpha$ for a
noncompact positive root $\alpha$, is $\pi_\lambda$.
\end{enumerate}

The semisimple versions apply to $U(p,q)$ through the
finite central covering $SU(p,q)\times U(1)\to U(p,q)$.
An irreducible unitary representation has a scalar
central character. Its restriction to $SU(p,q)$ is
irreducible, since a subspace invariant under that group
is then invariant under the whole group.
The character on the compact center has absolute value one;
probability measure on that center changes no
matrix-coefficient norm. The possible fixed measure factor
in the covering is absorbed in $b_\mu$ in (FD13).
The compact-type and root statements restrict to the
trace-zero hyperplane; all root differences have sum zero,
so adding the fixed central weight leaves every comparison
unchanged. The complexification of $SU(p,q)$ is $SL_n(\mathbb C)$,
which is simply connected, as required in the third input.

\subsection{An explicit scalar discrete series fixes the constant}

Let $s$ be an integer with $s>n-1$, and put
$\alpha=s-n>-1$. Consider the Hilbert space
\[
 \mathcal H_s=\left\{f\text{ holomorphic on }\mathcal B_{p,q}:
        \|f\|_s^2=
         \int_{\mathcal B_{p,q}}|f(Z)|^2
                         \det(I_q-Z^*Z)^{s-n}\,dZ<\infty\right\}.
 \tag{FD15}
\]
It is a closed subspace of the indicated $L^2$ space.
To see this, the mean-value estimate on a small ball about
any interior point, where the weight is bounded below,
controls point evaluation uniformly on compact subsets.
An $L^2$-Cauchy sequence of holomorphic functions therefore
converges locally uniformly to a holomorphic function,
which represents its $L^2$ limit. In particular every
point evaluation is continuous.

For $g\in G$ define $j(g,Z)=\det(CZ+D)$, and set
\[
       (\pi_s(g)f)(Z)=j(g^{-1},Z)^{-s}f(g^{-1}\cdot Z).
 \tag{FD16}
\]
The exponent is integral, so no choice of branch is involved.

\begin{lemma}\label{fd-scalar-model}
The representation $\pi_s$ is irreducible, unitary and
belongs to the discrete series. If
$I_s=I_{p,q}(s-n)$ is the integral in
\eqref{fd-ball-integral}, then
\[
                    d_{\mu_0}(\pi_s)=I_s^{-1}
                      =\pi^{-pq}
                        \prod_{i=1}^p\prod_{j=1}^q(s-i-j+1).
 \tag{FD17}\label{fd-scalar-degree}
\]
\end{lemma}
\begin{proof}
The identity
$j(gh,Z)=j(g,h\cdot Z)j(h,Z)$ proves the representation
law. Equations (FD6) and the Jacobian calculation prove
unitarity after a change of variables in (FD15).
The representation is strongly continuous: on a compact
subset of the ball the transformed holomorphic function
converges uniformly as $g\to1$. In the inner product with
a fixed $f$, the integral over the complementary tail is
uniformly bounded by its $L^2$ norm times $\|f\|_s$.
Letting that tail tend to zero proves
$\langle\pi_s(g)f,f\rangle\to\|f\|_s^2$, and unitarity
then gives norm convergence.

The constant function $1$ has norm squared $I_s$.
Averaging over the scalar rotations $Z\mapsto e^{it}Z$
and using the one-variable Cauchy formula gives
\[
                         \langle f,1\rangle_s=I_s f(0),
 \tag{FD18}
\]
with the inner product linear in the first argument.
Fubini is legitimate because $I_s<\infty$ and
$\int|f|\,\det(I-Z^*Z)^{s-n}\le\sqrt{I_s}\|f\|_s$.
Put $v=I_s^{-1/2}1$. For $Z=g\cdot0$, (FD16)--(FD18) give
\[
 |\langle f,\pi_s(g)v\rangle_s|^2
       =I_s|j(g,0)|^{-2s}|f(Z)|^2
       =I_s\det(I_q-Z^*Z)^s|f(Z)|^2.
 \tag{FD19}
\]
Thus the translates of $v$ span a dense subspace: a function
orthogonal to all of them vanishes at every point.

For $k=\operatorname{diag}(A,D)\in K$, the constant
has character $\det(D)^s$. The subspace with this
character is one-dimensional. Indeed its functions
are invariant under $Z\mapsto e^{it}Z$, and their
holomorphic power series have no terms of positive degree.
An orthogonal projection onto a closed invariant
subspace commutes with $K$, hence sends $v$ to either
$0$ or $v$. It then has the same alternative on every
translate of $v$, proving irreducibility by density.

Integrating (FD19) with (FD3)--(FD4) gives
\[
               \int_G|\langle f,\pi_s(g)v\rangle_s|^2\,d\mu_0(g)
                                  =I_s\|f\|_s^2.
 \tag{FD20}
\]
Consequently $f\mapsto I_s^{-1/2}
\langle f,\pi_s(\,\cdot\,)v\rangle_s$ is an isometric
intertwining map into the left regular representation,
with closed image. This proves the discrete-series claim.
Taking $f=v$ in (FD20) and comparing with (FD12) gives
$d_{\mu_0}(\pi_s)=I_s^{-1}$.
Finally, reversal of the factors in
\eqref{fd-ball-integral} gives
\[
 I_s=\pi^{pq}\prod_{j=0}^{q-1}
                       \frac{\Gamma(s-p-j)}{\Gamma(s-j)}.
\]
The gamma recurrence, obtained by integration by parts,
now proves (FD17).
\end{proof}

Its Harish--Chandra parameter is also explicit. In the
complexified Lie algebra with upper triangular positive
root spaces, differentiating (FD16) shows that the
constant vector is killed by every positive root vector
and has highest weight
\[
                 \Lambda=(\underbrace{0,\ldots,0}_{p},
                            \underbrace{s,\ldots,s}_{q}).
 \tag{FD21}
\]
For a cross-block upper root vector the derivative is a
negative coordinate derivative, so it kills a constant;
compact upper root vectors kill it as well.
The diagonal calculation gives $0$ in the first block
and $s$ in the second. The usual highest-weight
infinitesimal character is $\Lambda+\rho$, where
$\rho_i=(n+1)/2-i$. This follows from the
Harish--Chandra isomorphism, or from the ordered Capelli
eigenvalue calculation on a highest vector.
Only the unordered tuple is needed here.
Its within-block Vandermonde factors are $s_p,s_q$, and
its cross-block factors have absolute values
$s-p-j+i$ for $1\le i\le p$, $1\le j\le q$.
All are positive because $s>n-1$. Hence
\[
 \prod_{i<j}|(\Lambda+\rho)_i-(\Lambda+\rho)_j|
       =s_ps_q\prod_{i=1}^p\prod_{j=1}^q(s-i-j+1).
 \tag{FD22}
\]

\begin{theorem}\label{fd-formula}
For every discrete series of $U(p,q)$ with
Harish--Chandra parameter $\lambda$,
\begin{equation}\label{fd-exact}
            d_{\mu_{\rm cd}}(\pi_\lambda)
                 =\frac{\prod_{i<j}|\lambda_i-\lambda_j|}{s_n}.
\end{equation}
No extra lower bound on the noncompact root gaps is needed
in this classical argument.
\end{theorem}
\begin{proof}
Comparing (FD17), (FD22) and the parameter-independent
constant in (FD13) gives
$b_{\mu_0}=\pi^{-pq}/(s_ps_q)$.
Divide by $c_{p,q}$ using (FD10) and the scaling
rule following (FD12). The result is $b_{\mu_{\rm cd}}=1/s_n$.
\end{proof}

\paragraph{Compact signatures.}
The same formula holds when $pq=0$. In this case the group
is $U(n)$, and its compact-dual normalized Haar measure is
its probability measure. An irreducible unitary representation
of a compact group is finite dimensional. Schur orthogonality
for probability measure makes its formal degree its dimension.
If $\mu$ is its dominant integral highest weight and
$\lambda=\mu+\rho$, Weyl's dimension formula gives
\[
 \dim V_\mu
   =\prod_{i<j}\frac{\mu_i-\mu_j+j-i}{j-i}
   =\frac{\prod_{i<j}(\lambda_i-\lambda_j)}{s_n}.
\]
The equality $\prod_{i<j}(j-i)=s_n$ follows by grouping
the factors with fixed $j$. This is also the dimension
formula proved in Section~\ref{cd-section}.
For $n=0$ the group is trivial, its irreducible representation
is one dimensional, and both the empty Vandermonde and
$s_0$ are one. Thus its formal degree is one.

\subsection{Checking the source's compact-type definition}

Now allow the signs $\varepsilon_i\in\{1,-1\}$ to be in
any order, and suppose
\[
             \lambda_1>\cdots>\lambda_n,\qquad
             \lambda_i-(n-1)/2\in\mathbb Z.
 \tag{FD23}
\]
Positive roots are $e_i-e_j$ for $i<j$. A positive root
is compact if $\varepsilon_i=\varepsilon_j$ and
noncompact otherwise. Use their half sums
$\rho_{\rm c},\rho_{\rm n}$ and put
\[
 \rho=\rho_{\rm c}+\rho_{\rm n},\quad
 \rho_i=(n+1)/2-i,\quad
 \mu_0=\lambda+\rho_{\rm n}-\rho_{\rm c}.
 \tag{FD24}
\]
The compact Weyl group permutes the indices of each sign.
A compact highest weight $\mu$ is dominant when
$\mu_i\ge\mu_j$ for $i<j$ of the same sign.
The squared norm used in the source is
$\|\mu+2\rho_{\rm c}\|^2=\sum_i(\mu_i+2\rho_{{\rm c},i})^2$.

\begin{proposition}\label{fd-lowest-type}
The classical representation $\pi_\lambda$ contains the
compact type $\mu_0$ with multiplicity one. Every other
compact highest weight $\mu$ occurring in it satisfies
\[
             \|\mu+2\rho_{\rm c}\|^2
                        >\|\mu_0+2\rho_{\rm c}\|^2.
 \tag{FD25}
\]
Conversely, an irreducible unitary representation containing
the type $\mu_0$ and satisfying the corresponding weak
minimum inequality for all its compact types is
$\pi_\lambda$.
\end{proposition}
\begin{proof}
If $pq=0$, then $G=K$ and $\mu_0=\lambda-\rho$.
The corresponding compact representation has just this
one compact type, and any irreducible representation
containing it is that representation. This includes the
trivial group when $n=0$. We may therefore assume $pq>0$.
The strict integral gaps imply $\lambda_i-\lambda_j\ge j-i$.
For indices of the same sign let $b_{ij}$ count the
opposite-sign indices strictly between them. Directly
from the half sums,
\[
 \rho_{{\rm c},i}-\rho_{{\rm c},j}=j-i-b_{ij},\qquad
 \rho_{{\rm n},i}-\rho_{{\rm n},j}=b_{ij}.
\]
Thus
$\mu_{0,i}-\mu_{0,j}=\lambda_i-\lambda_j-(j-i)+2b_{ij}\ge0$.
Also $\mu_0=\lambda+\rho-2\rho_{\rm c}$ is integral:
$\lambda_i+\rho_i\in\mathbb Z$ and
$2\rho_{{\rm c},i}\in\mathbb Z$.
The tuple $\xi=\lambda+\rho_{\rm n}=\mu_0+\rho_{\rm c}$
is strictly decreasing within each sign.

For two tuples strictly decreasing within each sign, their
scalar product strictly decreases when one tuple is
permuted by a nonidentity element of $W_K$.
This is the rearrangement inequality: swapping an inverted
pair changes the scalar product by
$(a_i-a_j)(b_i-b_j)>0$, and repeated swaps sort each block.
Consequently, for $w\ne1$,
\[
                       \langle\lambda,w\xi-\xi\rangle<0.
\]
Every nonzero sum of positive noncompact roots instead
has positive scalar product with $\lambda$.
In (FD14) at $\mu=\mu_0$, only the identity permutation
can therefore contribute, and it contributes
$P_{\rm n}(0)=1$. This proves occurrence and multiplicity.

If $m_\mu(\pi_\lambda)>0$, some term of (FD14) is nonzero.
For its permutation $w$, write
\[
          w(\mu+\rho_{\rm c})=\lambda+\rho_{\rm n}+\beta,
 \qquad \beta\in\sum_{\alpha>0,\ {\rm noncompact}}\mathbb N\alpha.
\]
Put $x=\mu+\rho_{\rm c}$. The weak rearrangement
inequality applied within each sign gives
\begin{align*}
 \|\mu+2\rho_{\rm c}\|^2
   &=\|x\|^2+2\langle x,\rho_{\rm c}\rangle+\|\rho_{\rm c}\|^2\\
   &\ge\|wx+\rho_{\rm c}\|^2
     =\|\lambda+\rho+\beta\|^2
     \ge\|\lambda+\rho\|^2.
\end{align*}
The final inequality is strict if $\beta\ne0$, since
$\lambda+\rho$ has positive scalar product with every
positive root. If $\beta=0$, the two compact-dominant
tuples $x,\xi$ in the same $W_K$-orbit are equal, so
$\mu=\mu_0$. This proves (FD25).

For the converse, put $\nu=\lambda-\rho$. It is dominant
integral by (FD23). For a positive noncompact root
$\alpha=e_i-e_j$,
\[
 \|\mu_0-\alpha+2\rho_{\rm c}\|^2
       -\|\mu_0+2\rho_{\rm c}\|^2
       =-2\langle\lambda+\rho,\alpha\rangle+2<0.
 \tag{FD26}
\]
Indeed $\lambda_i-\lambda_j\ge j-i\ge1$, and
$\rho_i-\rho_j=j-i\ge1$.
Whenever $\mu_0-\alpha$ is compact dominant, its occurrence
would contradict the assumed minimum. All the hypotheses
of the stated compact-type characterization apply, with
$\nu+2\rho_{\rm n}=\mu_0$. It follows that the representation
is $\pi_\lambda$.
\end{proof}

\begin{corollary}\label{fd-source}
Suppose the source's Haar and representation data satisfy
their stated identifications: its measure is compact-dual
normalized; its representation is nonzero, irreducible,
unitary, strongly continuous and square integrable; and
it contains the compact type $\mu_0$ of (FD24), with
minimum squared norm as above. Then its formal degree
$d_\lambda$ satisfies
\[
                  s_n d_\lambda
                      =\prod_{i<j}(\lambda_i-\lambda_j).
 \tag{FD27}
\]
In particular the formal-degree assumption
\eqref{as-formal-degree-assumption} holds with equality
before its factor $1/2$ is inserted.
\end{corollary}
\begin{proof}
The source's equal-density measure conditions are exactly
(FD8)--(FD10), so Proposition~\ref{fd-measure-match}
identifies its Haar measure. Its compact-type conditions
are exactly those in Proposition~\ref{fd-lowest-type}.
Apply Theorem~\ref{fd-formula}; strict decrease removes
the absolute values.

There is also a direct way to apply the source's numerical
identification of its degree. That identification quantifies
over every representation with the specified properties
and every positive formal degree satisfying Schur
orthogonality. The classical $\pi_\lambda$ exists and
has those properties by (FD14) and
Proposition~\ref{fd-lowest-type}. Substituting this one
representation gives the same value for the stipulated
degree. Thus no inference from an empty class of
representations is used.
\end{proof}

\subsection{Normalization checks and scope}

For $p=q=1$, one has $c_{1,1}=1/\pi$, and the scalar
parameter difference is $s-1$. Equations (FD17) and
\eqref{fd-exact} give
\[
               d_{\mu_0}=(s-1)/\pi,\qquad
               d_{\mu_{\rm cd}}=s-1.
\]
More generally the constant is $s_n/(\pi^{pq}s_ps_q)$,
not just $\pi^{-pq}$. A claim that (FD27) holds for an
arbitrary Haar measure is false: replacing the measure
by twice itself divides its formal degree by two.
The source explicitly fixes the measure and is not
subject to this objection.

The compact-type condition also has real content.
For $U(1,1)$ the two discrete series in a regular
infinitesimal-character packet have different smallest
torus weights. Replacing the stipulated type by an
arbitrary occurring type would not justify the
parameter identification above. The proof retains
both its occurrence and its minimum property.

This section supplies the precise formal-degree value
and its normalization from the three named classical
theorems. It does not reprove those deep theorems or
claim a line-by-line reconstruction of the source's
independent discrete-series construction. It closes
the formal-degree input to the elementary size estimates.
The global multiplicity bound, the integrated Jacquet
identities and the arithmetic theta moments remain
separate questions.


% Rational orbital estimates and the complex multiplicity comparison.
% Labels in this fragment have the prefix amb-.
